\documentclass[10pt]{article}
% Shared LaTeX style for MIT 18.06 exam revision notes.
% Compile topic files with Tectonic, XeLaTeX, or LuaLaTeX.

\usepackage[a4paper,margin=0.72in]{geometry}
\usepackage{fontspec}
\usepackage{amsmath,amssymb,mathtools}
\usepackage{booktabs,array,tabularx}
\usepackage{enumitem}
\usepackage{titlesec}
\usepackage{fancyhdr}
\usepackage{microtype}
\usepackage{xcolor}

\setmainfont{Times New Roman}
\setsansfont{Arial}

\setlength{\parindent}{0pt}
\setlength{\parskip}{3.6pt}
\setlength{\abovedisplayskip}{6pt}
\setlength{\belowdisplayskip}{6pt}
\setlength{\abovedisplayshortskip}{4pt}
\setlength{\belowdisplayshortskip}{4pt}
\renewcommand{\arraystretch}{1.08}
\setlength{\tabcolsep}{4.5pt}

\titleformat{\section}
  {\normalfont\large\bfseries}
  {\thesection}
  {0.55em}
  {}
\titlespacing*{\section}{0pt}{10pt}{3pt}

\titleformat{\subsection}
  {\normalfont\normalsize\bfseries}
  {\thesubsection}
  {0.5em}
  {}
\titlespacing*{\subsection}{0pt}{7pt}{2pt}

\setlist[itemize]{leftmargin=1.35em,itemsep=1pt,topsep=2pt,parsep=0pt}
\setlist[enumerate]{leftmargin=1.55em,itemsep=1pt,topsep=2pt,parsep=0pt}

\pagestyle{fancy}
\fancyhf{}
\fancyfoot[C]{\scriptsize\color{black!45}MIT 18.06 Linear Algebra Exam Notes --- \thepage}
\renewcommand{\headrulewidth}{0pt}
\renewcommand{\footrulewidth}{0pt}

\newcommand{\notesubtitle}[1]{%
  {\centering\itshape #1\par}
  \vspace{2pt}
}

\newcommand{\source}[1]{%
  {\centering\small #1\par}
  \vspace{6pt}
}

\newcommand{\labelnote}[2]{%
  \par\smallskip
  \noindent\textbf{#1.}\quad #2
  \par\smallskip
}

\newcommand{\tightlabel}[1]{%
  \par\smallskip
  \noindent\textbf{#1.}\quad
}

\newcolumntype{Y}{>{\raggedright\arraybackslash}X}
\newcolumntype{L}[1]{>{\raggedright\arraybackslash}p{#1}}

\newenvironment{notetable}
  {\small\begin{tabularx}{\linewidth}}
  {\end{tabularx}}


\begin{document}

\begin{center}
{\LARGE Eigenvalue Quick Reference\par}
\end{center}
\notesubtitle{Concise revision notes for eigenvalue equations, checks, and basic matrix types.}
\source{Based on handwritten MIT 18.06 revision notes.}

\section{Core Idea}

\labelnote{Core idea}{Eigenvectors are directions that survive multiplication by \(A\). Along an eigenvector, the matrix only multiplies by a scalar.}

The defining equation is
\[
Ax=\lambda x,
\qquad x\ne0.
\]
Move all terms to one side:
\[
(A-\lambda I)x=0.
\]
For a nonzero solution, the matrix \(A-\lambda I\) must be singular:
\[
\det(A-\lambda I)=0.
\]

\section{Characteristic Equation}

For a \(2\times2\) matrix, the characteristic equation can be written as
\[
\lambda^2-\operatorname{trace}(A)\lambda+\det A=0.
\]
This is often the fastest way to check arithmetic.

\begin{center}
\small
\begin{tabularx}{\linewidth}{@{}L{0.26\linewidth}Y@{}}
\toprule
Quantity & Eigenvalue check \\
\midrule
Trace & Sum of eigenvalues. \\
Determinant & Product of eigenvalues. \\
Zero eigenvalue & Matrix is singular. \\
\bottomrule
\end{tabularx}
\end{center}

\section{Basic Examples}

\begin{center}
\small
\begin{tabularx}{\linewidth}{@{}L{0.32\linewidth}Y@{}}
\toprule
Matrix type & Eigenvalue memory \\
\midrule
Projection \(P\) with \(P^2=P\) & Eigenvalues are \(0\) or \(1\). \\
Permutation swapping two coordinates & Eigenvalues often include \(1\) and \(-1\). \\
Plane rotation by \(90^\circ\) & Eigenvalues are complex: \(i\) and \(-i\). \\
\bottomrule
\end{tabularx}
\end{center}

\section{Growth and Stability}

For repeated multiplication \(A^k\):
\[
|\lambda|<1
\]
means decay in that eigenvector direction.

\[
|\lambda|>1
\]
means growth in that direction.

\[
\lambda=1
\]
means that direction stays fixed.

\section{Common Mistakes}

\begin{center}
\small
\begin{tabularx}{\linewidth}{@{}L{0.40\linewidth}Y@{}}
\toprule
Mistake & Correct exam habit \\
\midrule
Allowing \(x=0\) as an eigenvector. & Eigenvectors must be nonzero. \\
Solving \(Ax=0\) for every eigenvector. & Solve \((A-\lambda I)x=0\). \\
Forgetting trace and determinant checks. & Use them after solving for \(\lambda\). \\
Assuming all rotations have real eigenvectors. & Plane rotations can have complex eigenvalues. \\
\bottomrule
\end{tabularx}
\end{center}

\section{Final Exam Checklist}

\tightlabel{Checklist}
\begin{enumerate}
  \item Write \(Ax=\lambda x\).
  \item Form \(A-\lambda I\).
  \item Solve \(\det(A-\lambda I)=0\).
  \item Check trace and determinant.
  \item For each \(\lambda\), solve \((A-\lambda I)x=0\).
  \item Use \(|\lambda|\) to judge growth or decay for powers.
\end{enumerate}

\end{document}
